% ============================================================
% style/packages.tex
%
% Пакеты документа.
%
% Назначение файла:
%   - подключить базовые пакеты;
%   - подключить пакеты для математики, графики, таблиц и рамок;
%   - подключить служебные пакеты для пользовательских команд;
%   - настроить TikZ / PGFPlots;
%   - подключить hyperref в конце.
%
% Важно:
%   этот вариант рассчитан на pdfLaTeX.
%
% Если позже нужен Times New Roman через XeLaTeX / LuaLaTeX,
% блок с fontenc/inputenc/babel нужно будет заменить на fontspec/polyglossia
% или fontspec/babel.
% ============================================================


% ------------------------------------------------------------
% 1. Русский и английский текст для pdfLaTeX
% ------------------------------------------------------------
%
% cmap      — корректный поиск и копирование кириллицы из PDF;
% fontenc   — кодировка шрифтов для кириллицы;
% inputenc  — входная кодировка исходного .tex-файла;
% babel     — переносы и языковые правила.
%
% Последний язык в babel становится основным.
% Здесь основной язык — русский.

\usepackage{cmap}
\usepackage[T2A]{fontenc}
\usepackage[utf8]{inputenc}
\usepackage[english,russian]{babel}


% ------------------------------------------------------------
% 2. Математика
% ------------------------------------------------------------

\usepackage{amsmath}
\usepackage{amssymb}
\usepackage{braket}
\usepackage{bm}


% ------------------------------------------------------------
% 3. Графика и изображения
% ------------------------------------------------------------
%
% graphicx  — вставка изображений;
% float     — управление плавающими объектами;
% wrapfig   — обтекание рисунков текстом;
% capt-of   — подписи вне стандартных окружений figure/table.

\usepackage{graphicx}
\usepackage{float}
\usepackage{wrapfig}
\usepackage{capt-of}


% ------------------------------------------------------------
% 4. Подписи к рисункам
% ------------------------------------------------------------
%
% labelformat=empty убирает стандартное "Рис. N".
% Это удобно, если подписи оформляются вручную или декоративно.
%
% subcaption нужен для подрисунков.

\usepackage[labelformat=empty]{caption}
\usepackage{subcaption}


% ------------------------------------------------------------
% 5. Таблицы
% ------------------------------------------------------------

\usepackage{array}
\usepackage{tabularx}
\usepackage{makecell}
\usepackage{multirow}
\usepackage{colortbl}


% ------------------------------------------------------------
% 6. Единицы измерения и специальные записи
% ------------------------------------------------------------

\usepackage{siunitx}
\usepackage{textgreek}
\usepackage{polynom}


% ------------------------------------------------------------
% 7. Списки
% ------------------------------------------------------------

\usepackage{enumitem}


% ------------------------------------------------------------
% 8. TikZ, PGFPlots и геометрические рисунки
% ------------------------------------------------------------
%
% tikz          — основной пакет для векторной графики;
% pgfplots      — графики функций и данных;
% tkz-euclide   — евклидова геометрия на TikZ;
% tikz-3dplot   — 3D-повороты и 3D-сцены.

\usepackage{tikz}
\usepackage{pgfplots}
\usepackage{tkz-euclide}
\usepackage{tikz-3dplot}


% ------------------------------------------------------------
% 9. PSTricks / pst-eucl
% ------------------------------------------------------------
%
% ВАЖНО:
%   pst-eucl относится к PSTricks.
%   Он может плохо работать напрямую с pdfLaTeX.
%
% Если при компиляции появятся ошибки, связанные с PSTricks,
% эту строку лучше закомментировать и заменить нужные рисунки на TikZ.

\usepackage{pst-eucl}


% ------------------------------------------------------------
% 10. Рамки, цветные блоки и декоративные вставки
% ------------------------------------------------------------
%
% tcolorbox — основа для rbox, rfig, rformula и других блоков;
% mdframed  — дополнительный пакет для рамок, если понадобится;
% darkmode  — тёмный режим страницы.

\usepackage[most]{tcolorbox}
\tcbuselibrary{xparse}
\usepackage{mdframed}
\usepackage{darkmode}


% ------------------------------------------------------------
% 11. Химические формулы
% ------------------------------------------------------------

\usepackage{chemfig}


% ------------------------------------------------------------
% 12. Оформление страницы и текста
% ------------------------------------------------------------
%
% fullpage  — уменьшает стандартные поля;
% changepage — локальное изменение ширины текста;
% ulem      — подчёркивания, например \uline;
% titlesec  — настройка \part, \chapter, \section и т.д.

\usepackage{fullpage}
\usepackage{changepage}
\usepackage{ulem}
\usepackage{titlesec}


% ------------------------------------------------------------
% 13. Служебные пакеты для пользовательских команд
% ------------------------------------------------------------
%
% xparse   — удобное создание команд через \NewDocumentCommand;
% etoolbox — проверки строк, например \ifstrempty и \ifstrequal;
% ifthen   — простая логика;
% calc     — вычисления с длинами.

\usepackage{xparse}
\usepackage{etoolbox}
\usepackage{ifthen}
\usepackage{calc}


% ------------------------------------------------------------
% 14. Библиотеки PGFPlots
% ------------------------------------------------------------

\usepgfplotslibrary{groupplots}
\usepgfplotslibrary{colormaps,patchplots}
\usepgfplotslibrary{fillbetween}


% ------------------------------------------------------------
% 15. Библиотеки TikZ
% ------------------------------------------------------------

\usetikzlibrary{arrows.meta}
\usetikzlibrary{calc,patterns,angles,quotes}
\usetikzlibrary{shadings}
\usetikzlibrary{decorations.markings,bending}
\usetikzlibrary{3d,backgrounds,intersections}
\usetikzlibrary{perspective}
\usetikzlibrary{shapes.arrows}


% ------------------------------------------------------------
% 16. Настройки графики
% ------------------------------------------------------------
%
% PointName=none и PointSymbol=none — настройки для pst-eucl.
%
% compat=1.16 оставлен как в исходном стиле.
% Если TeX Live свежий, можно будет поднять до более новой версии.

\psset{
    PointName=none,
    PointSymbol=none
}

\pgfplotsset{
    compat=1.16
}


% ------------------------------------------------------------
% 17. Тёмный режим
% ------------------------------------------------------------
%
% Делает страницу тёмной.
% Основные цвета самой темы задаются отдельно в:
%   style/colors.tex

\enabledarkmode


% ------------------------------------------------------------
% 18. Шрифт
% ------------------------------------------------------------
%
% В ChemNotes2 использовался mathpazo.
% В Physics Updated он не использовался.
%
% Сейчас оставляем стандартный Computer Modern.
%
% Если нужен вид ближе к ChemNotes2, можно раскомментировать:
%
% \usepackage{mathpazo}


% ------------------------------------------------------------
% 19. Ссылки
% ------------------------------------------------------------
%
% hyperref лучше подключать ближе к концу преамбулы.
%
% Цвета ссылок оставлены белыми, потому что colors.tex обычно
% подключается после packages.tex, и смысловые цвета rtext/rbg
% здесь ещё могут быть не определены.

\usepackage[
    unicode=true,
    colorlinks=true,
    linkcolor=white,
    urlcolor=white,
    citecolor=white,
    pdfborder={0 0 0}
]{hyperref}